\scanpage{14}
\mspage{4}

\[
\text{a)}\quad \varphi_!(e_{\widehat A})\longrightarrow e_{\mathcal X}\quad\text{is an epimorphism}
\]
\[
\text{b)}\quad \varphi_!(\alpha\times\beta)\longrightarrow
\varphi_!(\alpha)\times\varphi_!(\beta)\quad\text{is an epimorphism}
\qquad \forall\alpha,\beta\text{ in }\operatorname{Ob}A
\]
\[
\text{c)}\quad \varphi_!(\alpha\times_{\alpha_0}\beta)\longrightarrow
\varphi_!(\alpha)\times_{\varphi_!(\alpha_0)}\varphi_!(\beta)\quad\text{is an epimorphism}
\]
$\forall$ diagram
\diagram{\begin{tikzcd}[ampersand replacement=\&,column sep=small,row sep=small]
\alpha\arrow[dr]\&\&\beta\arrow[dl]\\
\&\alpha_0\&
\end{tikzcd}}
in $A$.

A particularly pleasant case where these conditions are satisfied is when $\varphi_!$ is left exact, i.e. that it corresponds \deleted{\unsure{to}} to a morphism of topoi
\[
f:\mathcal X\longrightarrow\widehat A
\]
by the formula
\[
\varphi_!=f^*.
\]
Then in a) b) c) the epimorphic arrows in question are even isomorphisms, and I presume that the converse is true as well. Yes -- and it is essentially \unsure{evident|evident}, on writing canonically, for an $F\in\operatorname{Ob}\widehat A$, $F=\varinjlim_{A/F}\alpha$ \unsure{$\ldots$}\unsure{.} One sees

\deleted{But \unsure{can one have}}
\inserted{in the same way for} b) and c)

\margin{13}
